% ============================================================================
% Capítulo 2: Probabilidade para Finanças — Revisão Rigorosa
% Módulo: Geometric Arbitrage Theory (DCA)
% ============================================================================

\chapter{Probabilidade para Finanças\texorpdfstring{\\}{:} Revisão Rigorosa}

A teoria moderna de finanças matemáticas repousa sobre os alicerces da teoria da probabilidade. O movimento browniano que modela a difusão de preços, os martingais que caracterizam a ausência de arbitragem e as expectativas condicionais que precificam derivativos --- todos esses conceitos são, em última instância, construções da teoria da probabilidade axiomatizada por Andrey Kolmogorov em 1933. Este capítulo oferece uma revisão rigorosa, porém acessível, dos conceitos probabilísticos essenciais para o estudo da Teoria de Arbitragem Geométrica, com ênfase em exemplos financeiros que ilustram cada construção teórica.

% ---------------------------------------------------------------------------
\section{Espaços de Probabilidade}
% ---------------------------------------------------------------------------

A formalização da probabilidade como ramo da matemática pura deve-se a Kolmogorov, cuja monografia \textit{Grundbegriffe der Wahrscheinlichkeitsrechnung} (1933) estabeleceu a axiomática que permanece como padrão até hoje.\footnote{Para uma exposição completa da axiomática de Kolmogorov, veja \textcite{williams1991probability}, Capítulos 1--3. \textcite{billingsley2012probability} oferece um tratamento mais extenso com ênfase em teoria da medida.}

\begin{defi}[$\sigma$-álgebra]
Seja $\Omega$ um conjunto não vazio. Uma coleção $\$\\mathcal{F}\$ \subseteq 2^\Omega$ é uma \textbf{$\sigma$-álgebra} sobre $\Omega$ se:
\begin{enumerate}[label=(\roman*)]
    \item $\Omega \in \F$;
    \item $A \in \$\\mathcal{F}\$ \implies A^c \in \F$ (fechamento sob complementação);
    \item $\{A_n\}_{n \in \N} \subseteq \$\\mathcal{F}\$ \implies \bigcup_{n=1}^\infty A_n \in \F$ (fechamento sob união contável).
\end{enumerate}
O par $(\Omega, \F)$ é denominado \textbf{espaço mensurável}.
\end{defi}

Em finanças, uma $\sigma$-álgebra modela a \textbf{estrutura de informação} disponível em determinado instante. Considere o preço de uma ação observado diariamente: a $\sigma$-álgebra $\$\\mathcal{F}_t\$$ contém todos os eventos cuja ocorrência é conhecida até o dia $t$. Esta interpretação será fundamental quando introduzirmos filtrações na Seção~\ref{sec:martingais}.

\begin{defi}[Medida de Probabilidade]
Uma função $\P: \$\\mathcal{F}\$ \to [0,1]$ é uma \textbf{medida de probabilidade} se satisfaz os \textbf{axiomas de Kolmogorov}:
\begin{enumerate}[label=(K\arabic*)]
    \item $\P(\Omega) = 1$ (normalização);
    \item $\P(A) \geq 0$ para todo $A \in \F$ (não-negatividade);
    \item Para toda coleção $\{A_n\}_{n \in \N} \subseteq \F$ de conjuntos disjuntos dois a dois, vale a $\sigma$-aditividade:
    \begin{equation}
        \P\brk{\bigcup_{n=1}^\infty A_n} = \sum_{n=1}^\infty \P(A_n).
    \end{equation}
\end{enumerate}
A tripla $(\Omega, \F, \P)$ é um \textbf{espaço de probabilidade}.
\end{defi}

\begin{ex}[Lançamento de moeda infinita]
Seja $\Omega = \{0,1\}^\N$ o espaço de todas as sequências infinitas de caras (1) e coroas (0). A $\sigma$-álgebra gerada pelos cilindros finitos --- conjuntos da forma $\{\omega \in \Omega : \omega_1 = a_1, \dots, \omega_n = a_n\}$ --- é a $\sigma$-álgebra produto. A medida de probabilidade produto $\P = \bigotimes_{n=1}^\infty \mu_n$, com $\mu_n(\{0\}) = \mu_n(\{1\}) = 1/2$, define o espaço de probabilidade canônico para uma sequência de variáveis de Bernoulli independentes. Em finanças, este espaço modela o passeio aleatório simétrico, um dos alicerces do modelo binomial de precificação de opções.
\end{ex}

\begin{defi}[Filtração]
Uma família crescente de $\sigma$-álgebras $\{\$\\mathcal{F}_t\$\}_{t \geq 0}$ com $\F_s \subseteq \$\\mathcal{F}_t\$$ para $s \leq t$ é uma \textbf{filtração}. O espaço $(\Omega, \F, \{\$\\mathcal{F}_t\$\}, \P)$ é um \textbf{espaço de probabilidade filtrado}. A filtração aumentada usual\footnote{A filtração usual inclui as condições de completude e continuidade à direita, necessárias para a teoria geral de processos estocásticos. Veja \textcite{oksendal2003stochastic}, Capítulo 2.} satisfaz $\$\\mathcal{F}_t\$ = \bigcap_{s > t} \F_s$ (continuidade à direita) e $\F_0$ contém todos os conjuntos $\P$-nulos.
\end{defi}

A filtração $\$\\mathcal{F}_t\$$ representa a informação acumulada até o instante $t$. Em um mercado financeiro, $\$\\mathcal{F}_t\$$ contém a história completa dos preços $S_u$ para $0 \leq u \leq t$, mas não revela valores futuros. Esta assimetria informacional é o cerne da não-antecipatividade: uma estratégia de investimento $\theta_t$ deve ser $\$\\mathcal{F}_t\$$-mensurável, ou seja, baseada apenas em informação passada e presente.

% ---------------------------------------------------------------------------
\section{Variáveis Aleatórias e Esperança}
% ---------------------------------------------------------------------------

\begin{defi}[Variável Aleatória]
Seja $(\Omega, \F, \P)$ um espaço de probabilidade. Uma função $X: \Omega \to \R$ é uma \textbf{variável aleatória} (v.a.) se $X$ é $\F$-mensurável, isto é, $X^{-1}(B) \in \F$ para todo boreliano $B \in \mathcal{B}(\R)$. A \textbf{distribuição} de $X$ é a medida imagem $\P_X := \P \circ X^{-1}$ sobre $(\R, \mathcal{B}(\R))$.
\end{defi}

A condição de mensurabilidade garante que possamos atribuir probabilidades a eventos do tipo $\{X \leq a\}$, $\{X \in (a,b]\}$, etc. Em termos financeiros, o retorno logarítmico $R_t = \log(S_t/S_{t-1})$ é uma variável aleatória cuja distribuição determina o risco do ativo.

\begin{defi}[Esperança]
A \textbf{esperança} (ou valor esperado) de uma v.a. $X \geq 0$ é definida como a integral de Lebesgue:
\begin{equation}
    \E[X] := \int_\Omega X(\omega) \dif \P(\omega).
\end{equation}
Para $X$ qualquer, decompomos $X = X^+ - X^-$ (parte positiva e negativa) e definimos $\E[X] := \E[X^+] - \E[X^-]$, desde que pelo menos uma das integrais seja finita. Dizemos que $X \in \Lp^1(\Omega, \F, \P)$ se $\E[|X|] < \infty$.
\end{defi}

\begin{defi}[Momentos e Variância]
Para $k \in \N$, o \textbf{$k$-ésimo momento} (se existir) é $\E[X^k]$. O \textbf{$k$-ésimo momento central} é $\E[(X - \E[X])^k]$. A \textbf{variância} é o segundo momento central:
\begin{equation}
    \Var(X) := \E[(X - \E[X])^2] = \E[X^2] - (\E[X])^2.
\end{equation}
A \textbf{covariância} entre duas v.a. $X, Y \in \Lp^2$ é $\Cov(X,Y) := \E[(X - \E[X])(Y - \E[Y])]$.
\end{defi}

Em finanças, a variância (ou sua raiz quadrada, a volatilidade) é a medida de risco mais elementar. A covariância entre retornos de ativos distintos é a base da teoria de carteiras de Markowitz e do Capital Asset Pricing Model (CAPM).

\begin{defi}[Função Característica]
A \textbf{função característica} de uma v.a. $X$ é $\phi_X(u) := \E[e^{iuX}]$, para $u \in \R$. Esta função \textbf{sempre existe} (pois $|e^{iuX}| = 1$), é uniformemente contínua e determina unicamente a distribuição de $X$\footnote{Teorema da unicidade de Lévy. \textcite{williams1991probability}, Seção 16.6.}.
\end{defi}

A função característica é uma ferramenta indispensável em finanças quantitativas, especialmente na precificação de opções via transformada de Fourier (método de Carr--Madan) e no estudo de distribuições infinitamente divisíveis que modelam retornos em tempo contínuo.

\begin{teo}[Desigualdade de Markov e Chebyshev]
Seja $X$ uma v.a. não-negativa. Para todo $\lambda > 0$:
\begin{equation}
    \P(X \geq \lambda) \leq \frac{\E[X]}{\lambda} \quad \text{(Markov)}.
\end{equation}
Se $X \in \Lp^2$, então para todo $\varepsilon > 0$:
\begin{equation}
    \P(|X - \E[X]| \geq \varepsilon) \leq \frac{\Var(X)}{\varepsilon^2} \quad \text{(Chebyshev)}.
\end{equation}
\end{teo}

Estas desigualdades fornecem cotas para probabilidades de cauda e são a base teórica para medidas de risco do tipo Value-at-Risk (VaR).

% ---------------------------------------------------------------------------
\section{Distribuições Relevantes}
% ---------------------------------------------------------------------------

Apresentamos as quatro distribuições de probabilidade mais importantes para modelagem financeira. Cada uma captura um aspecto distinto do comportamento de preços e riscos.

\begin{defi}[Distribuição Normal -- Gaussiana]
$X \sim \mathcal{N}(\mu, \sigma^2)$ se possui densidade
\begin{equation}
    f_X(x) = \frac{1}{\sigma\sqrt{2\pi}} \exp\brk{-\frac{(x-\mu)^2}{2\sigma^2}}, \quad x \in \R.
\end{equation}
com $\E[X] = \mu$ e $\Var(X) = \sigma^2$.
\end{defi}

No modelo de Black--Scholes, o log-retorno $\log(S_t/S_0)$ segue uma distribuição normal com média $(\mu - \sigma^2/2)t$ e variância $\sigma^2 t$. Embora a hipótese de normalidade seja frequentemente violada na prática (caudas pesadas e assimetria nos retornos reais), a distribuição normal permanece como ponto de partida fundamental.

\begin{defi}[Distribuição Log-Normal]
$X \sim \text{LogNormal}(\mu, \sigma^2)$ se $\log X \sim \mathcal{N}(\mu, \sigma^2)$. A densidade é
\begin{equation}
    f_X(x) = \frac{1}{x\sigma\sqrt{2\pi}} \exp\brk{-\frac{(\log x - \mu)^2}{2\sigma^2}}, \quad x > 0,
\end{equation}
com $\E[X] = e^{\mu + \sigma^2/2}$ e $\Var(X) = e^{2\mu + \sigma^2}(e^{\sigma^2} - 1)$.
\end{defi}

O preço de uma ação no modelo de Black--Scholes é log-normal: $S_t = S_0 \exp((\mu - \sigma^2/2)t + \sigma W_t)$, onde $W_t$ é um movimento browniano. A log-normalidade garante que $S_t > 0$ quase certamente, respeitando a responsabilidade limitada dos acionistas.

\begin{defi}[Distribuição de Poisson]
$X \sim \text{Poisson}(\lambda)$ se $\P(X = k) = e^{-\lambda} \lambda^k / k!$ para $k \in \N_0$, com $\E[X] = \Var(X) = \lambda$.
\end{defi}

Em finanças, a distribuição de Poisson modela a \textbf{chegada de eventos raros}: falências, defaults de crédito, saltos de preço, ou a chegada de ordens de compra e venda em mercados de alta frequência. Processos de Poisson compostos são a base dos modelos de difusão com saltos (Merton, 1976; Kou, 2002).

\begin{defi}[Distribuição Exponencial]
$X \sim \text{Exp}(\lambda)$ se possui densidade $f_X(x) = \lambda e^{-\lambda x}$ para $x \geq 0$, com $\E[X] = 1/\lambda$ e $\Var(X) = 1/\lambda^2$. A propriedade fundamental é a \textbf{ausência de memória}:
\begin{equation}
    \P(X > s + t \mid X > s) = \P(X > t), \quad \forall s, t \geq 0.
\end{equation}
\end{defi}

O tempo entre eventos consecutivos em um processo de Poisson é exponencial. Em modelagem de risco de crédito, o tempo até o default de uma firma é modelado como uma variável exponencial (ou sua generalização, a distribuição de Weibull). A ausência de memória implica que a \textit{taxa de hazard} (intensidade de default) é constante, uma hipótese simplificadora que pode ser relaxada com processos de Cox (Poisson duplamente estocásticos).

\begin{ex}[Retorno de carteira]
Considere uma carteira com $n$ ativos cujos retornos $R_i \sim \mathcal{N}(\mu_i, \sigma_i^2)$ são conjuntamente normais. O retorno da carteira $R_p = \sum_{i=1}^n w_i R_i$ (com pesos $w_i$) também é normal: $R_p \sim \mathcal{N}(\sum w_i \mu_i, \sum_{i,j} w_i w_j \sigma_{ij})$. Esta propriedade de fechamento sob combinações lineares explica a ubiquidade da distribuição normal em finanças, embora a hipótese de normalidade conjunta seja uma idealização.
\end{ex}

% ---------------------------------------------------------------------------
\section{Probabilidade Condicional e Esperança Condicional}
% ---------------------------------------------------------------------------

O conceito de esperança condicional é a ferramenta matemática mais importante da teoria de precificação de ativos. Diferentemente da probabilidade condicional elementar $\P(A \mid B) = \P(A \cap B)/\P(B)$, que condiciona em um evento de probabilidade positiva, a esperança condicional generaliza este conceito para condicionamento em $\sigma$-álgebras --- ou seja, em \textbf{conjuntos de informação}.

\begin{defi}[Esperança Condicional]
Seja $X \in \Lp^1(\Omega, \F, \P)$ e $\G \subseteq \F$ uma sub-$\sigma$-álgebra. A \textbf{esperança condicional} de $X$ dada $\G$ é a única (q.c.) variável aleatória $\E[X \mid \G]$ satisfazendo:
\begin{enumerate}[label=(\roman*)]
    \item $\E[X \mid \G]$ é $\G$-mensurável;
    \item Para todo $G \in \G$, $\int_G \E[X \mid \G] \dif \P = \int_G X \dif \P$.
\end{enumerate}
\end{defi}

A condição (ii) afirma que $\E[X \mid \G]$ é a \textbf{projeção} de $X$ no subespaço das funções $\G$-mensuráveis no sentido de $\Lp^2$, ou seja, é a melhor aproximação de $X$ usando apenas a informação contida em $\G$. Em finanças, se $\G = \$\\mathcal{F}_t\$$ representa a informação até $t$, então $\E[X \mid \$\\mathcal{F}_t\$]$ é a melhor previsão de $X$ com base em toda a história disponível.

\begin{teo}[Propriedades da Esperança Condicional]
Sejam $X, Y \in \Lp^1$ e $\G, \H$ sub-$\sigma$-álgebras de $\F$. Valem:
\begin{enumerate}[label=(\roman*)]
    \item \textbf{Linearidade:} $\E[aX + bY \mid \G] = a\E[X \mid \G] + b\E[Y \mid \G]$;
    \item \textbf{Retirada do conhecido:} se $Y$ é $\G$-mensurável e limitada, $\E[YX \mid \G] = Y \E[X \mid \G]$;
    \item \textbf{Independência:} se $X$ é independente de $\G$, então $\E[X \mid \G] = \E[X]$;
    \item \textbf{Propriedade de torre:} se $\H \subseteq \G$, então
    \begin{equation}
        \E\brk{\E[X \mid \G] \mid \H} = \E[X \mid \H];
    \end{equation}
    \item \textbf{Monotonicidade:} $X \leq Y$ q.c. $\implies$ $\E[X \mid \G] \leq \E[Y \mid \G]$ q.c.;
    \item \textbf{Desigualdade de Jensen condicional:} se $\phi$ é convexa e $\phi(X) \in \Lp^1$, então $\phi(\E[X \mid \G]) \leq \E[\phi(X) \mid \G]$ q.c.
\end{enumerate}
\end{teo}

A propriedade de torre (iv) é onipresente em finanças. Por exemplo, ao precificar uma opção americana por indução retroativa, a cada passo calculamos $\E[\text{payoff futuro} \mid \$\\mathcal{F}_t\$]$, e a propriedade de torre garante a consistência temporal do algoritmo de programação dinâmica.

\begin{ex}[Precificação neutra ao risco]
Em um mercado completo, o preço de um derivativo com payoff $H_T$ no vencimento $T$ é dado por
\begin{equation}
    V_t = e^{-r(T-t)} \, \E^\Q [H_T \mid \$\\mathcal{F}_t\$],
\end{equation}
onde $\Q$ é a medida martingal equivalente (medida neutra ao risco) e $r$ é a taxa livre de risco. A esperança condicional $\E^\Q[\cdot \mid \$\\mathcal{F}_t\$]$ representa o operador de precificação: ele calcula o valor esperado descontado do payoff futuro, condicionado a toda informação disponível no instante $t$. A propriedade de torre assegura que $e^{-rt} V_t$ é um $\Q$-martingal, implicando ausência de arbitragem.
\end{ex}

% ---------------------------------------------------------------------------
\section{Martingais}
\label{sec:martingais}
% ---------------------------------------------------------------------------

O conceito de martingal está no coração da teoria moderna de finanças. Intuitivamente, um martingal modela um \textbf{jogo justo}: a expectativa do valor futuro, condicionada à informação presente, é exatamente o valor presente. Esta propriedade é a tradução matemática precisa da \textbf{ausência de oportunidades de arbitragem} em mercados financeiros\footnote{O Primeiro Teorema Fundamental de Precificação de Ativos afirma que um mercado é livre de arbitragem se, e somente se, existe uma medida martingal equivalente. Veja \textcite{delbaen2006mathematics} e \textcite{williams1991probability}, Capítulo 10.}.

\begin{defi}[Martingal]
Seja $(\Omega, \F, \{\$\\mathcal{F}_t\$\}_{t \geq 0}, \P)$ um espaço de probabilidade filtrado. Um processo estocástico $\{M_t\}_{t \geq 0}$ adaptado a $\{\$\\mathcal{F}_t\$\}$ com $M_t \in \Lp^1$ para todo $t$ é um:
\begin{itemize}
    \item \textbf{martingal} se $\E[M_t \mid \F_s] = M_s$ para todo $0 \leq s \leq t$;
    \item \textbf{submartingal} se $\E[M_t \mid \F_s] \geq M_s$ para todo $0 \leq s \leq t$;
    \item \textbf{supermartingal} se $\E[M_t \mid \F_s] \leq M_s$ para todo $0 \leq s \leq t$.
\end{itemize}
\end{defi}

A distinção entre sub e supermartingal reflete a direção da tendência: um submartingal tende a crescer em média (jogo favorável ao apostador), enquanto um supermartingal tende a decrescer (jogo desfavorável).

\begin{ex}[Passeio Aleatório Simétrico]
Seja $\{X_n\}_{n \in \N}$ uma sequência de v.a. i.i.d. com $\P(X_n = 1) = \P(X_n = -1) = 1/2$. Defina $M_n = \sum_{k=1}^n X_k$ com $M_0 = 0$ e a filtração natural $\F_n = \sigma(X_1, \dots, X_n)$. Então $\{M_n\}$ é um martingal, pois
\begin{equation}
    \E[M_{n+1} \mid \F_n] = \E[M_n + X_{n+1} \mid \F_n] = M_n + \E[X_{n+1}] = M_n.
\end{equation}
O passeio aleatório é o análogo discreto do movimento browniano e constitui a base do modelo binomial de Cox--Ross--Rubinstein (CRR).
\end{ex}

\begin{ex}[Preço Descontado sob Medida Neutra ao Risco]
Seja $S_t$ o preço de um ativo e $B_t = e^{rt}$ o valor de um título livre de risco (conta bancária). Sob a medida martingal equivalente $\Q$, o \textbf{preço descontado} $\widetilde{S}_t := S_t / B_t = e^{-rt} S_t$ é um $\Q$-martingal:
\begin{equation}
    \E^\Q[\widetilde{S}_t \mid \F_s] = \widetilde{S}_s, \quad 0 \leq s \leq t.
\end{equation}
Esta é a condição central de não-arbitragem: o valor presente de qualquer ativo, descontado pela taxa livre de risco, deve ser um martingal sob $\Q$. No modelo de Black--Scholes, $\dif \widetilde{S}_t = \sigma \widetilde{S}_t \dif W_t^\Q$, onde $W_t^\Q$ é um $\Q$-movimento browniano, e a representação como integral estocástica de um martingal é explícita.
\end{ex}

\begin{teo}[Decomposição de Doob--Meyer]
Seja $\{X_t\}$ um submartingal contínuo à direita. Então existe uma decomposição única
\begin{equation}
    X_t = M_t + A_t,
\end{equation}
onde $\{M_t\}$ é um martingal e $\{A_t\}$ é um processo crescente previsível com $A_0 = 0$.
\end{teo}

Esta decomposição separa um processo estocástico em sua componente imprevisível (martingal) e sua tendência previsível (\textit{drift}). Em finanças, a decomposição de Doob--Meyer fundamenta a distinção entre risco sistemático (tendência) e risco idiossincrático (martingal), e é essencial para a fórmula de Itô e o cálculo estocástico.

\begin{teo}[Teorema do Sampling Opcional de Doob]
Seja $\{M_n\}_{n \geq 0}$ um martingal e $\tau$ um tempo de parada limitado. Então
\begin{equation}
    \E[M_\tau \mid \F_0] = M_0.
\end{equation}
Em particular, $\E[M_\tau] = \E[M_0]$.
\end{teo}

Este teorema justifica que, em um jogo justo, nenhuma estratégia de parada baseada apenas em informação passada (\textit{stopping time}) pode gerar lucro esperado positivo. Em finanças, isto implica que não se pode ``bater o mercado'' sistematicamente sem assumir risco adicional --- uma formalização da Hipótese do Mercado Eficiente (EMH) no contexto de martingais.

% ---------------------------------------------------------------------------
\section{Movimento Browniano}
% ---------------------------------------------------------------------------

O movimento browniano (ou processo de Wiener) é o objeto matemático central da modelagem financeira em tempo contínuo. Sua onipresença decorre do Teorema Central do Limite funcional (Teorema de Donsker): o movimento browniano é o limite de escala de passeios aleatórios, e portanto surge naturalmente como modelo para flutuações agregadas de preços resultantes de muitas decisões independentes de agentes de mercado.

\begin{defi}[Movimento Browniano Padrão]
Um processo estocástico $\{W_t\}_{t \geq 0}$ definido em $(\Omega, \F, \{\$\\mathcal{F}_t\$\}, \P)$ é um \textbf{movimento browniano padrão} (ou processo de Wiener) se:
\begin{enumerate}[label=(B\arabic*)]
    \item $W_0 = 0$ quase certamente;
    \item \textbf{Incrementos independentes:} para todo $0 \leq t_1 < t_2 < \cdots < t_n$, as v.a. $W_{t_2} - W_{t_1}, \dots, W_{t_n} - W_{t_{n-1}}$ são independentes;
    \item \textbf{Incrementos estacionários gaussianos:} para $0 \leq s < t$, $W_t - W_s \sim \mathcal{N}(0, t-s)$;
    \item \textbf{Trajetórias contínuas:} $t \mapsto W_t(\omega)$ é contínua para quase todo $\omega$.
\end{enumerate}
\end{defi}

As quatro propriedades caracterizam completamente o movimento browniano. A condição (B3) implica que $\E[W_t] = 0$ e $\E[W_t W_s] = \min(s,t)$ para todo $s,t \geq 0$.

\begin{teo}[Propriedades do Movimento Browniano]\label{teo:propriedades-mb}
Seja $\{W_t\}$ um movimento browniano padrão. Valem as seguintes propriedades:
\begin{enumerate}[label=(\roman*)]
    \item \textbf{Martingal:} $W_t$ é um $\$\\mathcal{F}_t\$$-martingal.
    \item \textbf{Variação quadrática:} $\langle W \rangle_t = t$ quase certamente, e a partição refinada satisfaz
    \begin{equation}
        \lim_{\|\Pi\| \to 0} \sum_{t_i \in \Pi} (W_{t_{i+1}} - W_{t_i})^2 = t \quad \text{em } \Lp^2.
    \end{equation}
    \item \textbf{Não-diferenciabilidade:} para quase todo $\omega$, a trajetória $t \mapsto W_t(\omega)$ não é diferenciável em nenhum ponto.
    \item \textbf{Propriedade de escala:} $\{c^{-1/2} W_{ct}\}_{t \geq 0}$ é um movimento browniano para todo $c > 0$.
    \item \textbf{Lei do logaritmo iterado:}
    \begin{equation}
        \limsup_{t \to 0^+} \frac{W_t}{\sqrt{2t \log\log(1/t)}} = 1 \quad \text{q.c.}
    \end{equation}
    \item \textbf{Propriedade forte de Markov:} $W_{t+\tau} - W_\tau$ é um movimento browniano independente de $\F_\tau$, para todo tempo de parada $\tau$ finito q.c.
\end{enumerate}
\end{teo}

A variação quadrática não-nula é a propriedade mais crucial para o cálculo estocástico. Intuitivamente, $(W_{t+\Delta} - W_t)^2 \approx \Delta$ em média, o que implica que $\dif W_t)^2 = \dif t$ na notação diferencial. Este fato é a origem do termo de segunda ordem na fórmula de Itô, que distingue o cálculo estocástico do cálculo usual.

\begin{defi}[Movimento Browniano Geométrico]
Um processo $\{S_t\}_{t \geq 0}$ é um \textbf{movimento browniano geométrico} (MBG) com drift $\mu$ e volatilidade $\sigma$ se satisfaz a equação diferencial estocástica (EDE)
\begin{equation}
    \frac{\dif S_t}{S_t} = \mu \dif t + \sigma \dif W_t,
\end{equation}
cuja solução é $S_t = S_0 \exp\brk{(\mu - \sigma^2/2)t + \sigma W_t}$.
\end{defi}

O MBG é o modelo-padrão para preços de ações no paradigma de Black--Scholes--Merton. A interpretação intuitiva é que o retorno relativo $\dif S_t / S_t$ tem uma componente determinística $\mu \dif t$ (tendência) e uma componente aleatória $\sigma \dif W_t$ (volatilidade). A aplicação do lema de Itô a $f(S_t) = \log S_t$ produz o fator de correção $-\sigma^2/2$, essencial para garantir que $\E[S_t] = S_0 e^{\mu t}$.

\begin{ex}[Construção de Lévy do Movimento Browniano]\label{ex:construcao-levy}
A existência do movimento browniano pode ser estabelecida construtivamente. Define-se $W_t$ para $t$ diádico por interpolação linear com incrementos gaussianos independentes: para cada nível $n$, os valores $W_{k/2^n}$ são definidos recursivamente. Para $t$ arbitrário, define-se $W_t$ como o limite uniforme das aproximações lineares por partes. A convergência uniforme quase certa é garantida pelo Lema de Borel--Cantelli. Esta construção, devida a Paul Lévy (1939), demonstra que o espaço $C[0,\infty)$ das funções contínuas (com a topologia da convergência uniforme em compactos) é o espaço amostral canônico do movimento browniano, com a medida de Wiener $\mathbb{W}$ sobre a $\sigma$-álgebra de Borel de $C[0,\infty)$.\footnote{Para detalhes completos da construção de Lévy e da medida de Wiener, consulte \textcite{oksendal2003stochastic}, Apêndice A, e \textcite{billingsley2012probability}, Seção 37.}
\end{ex}

% ---------------------------------------------------------------------------
\section*{Resumo do Capítulo}
% ---------------------------------------------------------------------------

Este capítulo estabeleceu o ferramental probabilístico necessário para o estudo da Teoria de Arbitragem Geométrica. Partindo dos axiomas de Kolmogorov, construímos progressivamente os conceitos de variável aleatória, esperança condicional, martingal e movimento browniano. Cada conceito foi ilustrado com exemplos financeiros concretos: a $\sigma$-álgebra como estrutura de informação, a esperança condicional como operador de precificação, o martingal como condição de não-arbitragem, e o movimento browniano como modelo-limite de flutuações de preços.

Os capítulos subsequentes utilizarão intensivamente este ferramental. O cálculo estocástico de Itô (Capítulo~3) estenderá o movimento browniano a integrais estocásticas, permitindo a modelagem de estratégias de investimento contínuas. A Teoria de Arbitragem Geométrica propriamente dita (Capítulo~4) aplicará martingais e medidas equivalentes para caracterizar geometricamente a ausência de arbitragem em mercados com fricções.

\begin{center}
\rule{0.5\textwidth}{0.4pt}
\end{center}

% ---------------------------------------------------------------------------
% Referências bibliográficas citadas neste capítulo
% ---------------------------------------------------------------------------
\begin{thebibliography}{9}

\bibitem{williams1991probability}
Williams, D.
\emph{Probability with Martingales}.
Cambridge University Press, 1991.
\url{https://doi.org/10.1017/CBO9780511813658}

\bibitem{billingsley2012probability}
Billingsley, P.
\emph{Probability and Measure}, Anniversary Edition.
Wiley, 2012.
ISBN: 978-1-118-12237-2.

\bibitem{oksendal2003stochastic}
\O{}ksendal, B.
\emph{Stochastic Differential Equations: An Introduction with Applications}, 6th ed.
Springer, 2003.
\url{https://doi.org/10.1007/978-3-642-14394-6}

\bibitem{delbaen2006mathematics}
Delbaen, F.; Schachermayer, W.
\emph{The Mathematics of Arbitrage}.
Springer, 2006.

\end{thebibliography}


